\clearpage\begingroup\NoSourceHead\SourceLeaf{538b}
\renewcommand{\thefootnote}{\arabic{footnote})}\setcounter{footnote}{8}
\newcommand{\BolzaEqualAboveGreater}{\mathrel{\vcenter{\offinterlineskip\hbox{$\scriptstyle=$}\kern-.15ex\hbox{$>$}}}}
\ArithmeticHead{81}{R}{EIN EXISTENZTHEOREM DER POTENTIALTHEORIE.}
\noindent\hspace*{\parindent}4. Nachdem so gezeigt ist, dass in den belegten Teilen der Fläche die
Funktion $W$ konstant sein muss, etwa
\begin{equation*}\tag{5}
W=A,
\end{equation*}
bleiben noch die etwa vorhandenen unbelegten Teile der Fläche zu betrachten.
In diesen Teilen kommt wegen (2) zu den bisherigen Bedingungen für $\mu$ noch
die weitere hinzu, dass $\mu\BolzaEqualAboveGreater0$ sein muss. Hieraus erhält \textsc{Gauss} das Resultat,
dass in den nicht belegten Teilen der Fläche
\begin{equation*}\tag{6}
W\BolzaEqualAboveGreater A
\end{equation*}
sein muss. Den Beweis hierfür überlässt er dem Leser; ein solcher ergibt
sich in der Tat auch ohne weiteres durch eine leichte Modifikation der früheren
Schlussweise. Wäre nämlich in einem noch so kleinen nicht belegten Teil $q$
der Fläche $W<A$, so sei $p$ eines der sicher vorhandenen belegten Stücke;
alsdann wähle man mit dieser neuen Bedeutung von $p$ und $q$ die Funktion $\mu$
genau so wie im vorigen Fall. Dann führt die Formel (4) zusammen mit der
in den belegten Teilen geltenden Gleichung (5) zu einem negativen Wert von
$\delta\Omega$, womit die Behauptung (6) bewiesen ist.
\par\noindent\hspace*{\parindent}Für die Variationsrechnung ist gerade dieser zweite Teil des \textsc{Gauss}schen
Resultates besonders interessant, weil hier wohl das erste Beispiel eines \textls[100]{Variationsproblems mit einer Ungleichung als Nebenbedingung} vorliegt, eine Klasse von Aufgaben, die später von \textsc{Weierstrass}\footnote{Vgl. \textsc{Kneser}, \textit{Lehrbuch der Variationsrechnung}, § 44; \textsc{Bolza}, \textit{Vorlesungen über Variationsrechnung},
§ 52 und \textsc{Hadamard}, \textit{Leçons sur le calcul des variations}, Nr. 161—164, 213 bis.} in seinen Vorlesungen behandelt worden ist.
\par\noindent\hspace*{\parindent}\textsc{Weierstrass} hat auch, und zwar aus der Betrachtung der ersten Variation,
die Bedingungen abgeleitet, die an den Übergangsstellen, in denen das Gleichheitszeichen in das Ungleichheitszeichen übergeht, erfüllt sein müssen, und die
zur Bestimmung dieser Stellen dienen. Man könnte versucht sein, durch ein
dem \textsc{Weierstrass}schen analoges Verfahren die Ausdehnung des nicht belegten
Teiles der Fläche bestimmen zu wollen. Dass man aber auf diese Weise nicht
zum Ziele kommen würde, geht a priori daraus hervor, dass aus den Be\mbox{-}
\par\noindent\hbox to\linewidth{\small X 2 Abh. 5.\hfill 11\hspace*{5em}}\par
\clearpage\SourceLeaf{539}
\ArithmeticHead{82}{L}{BOLZA, GAUSS UND DIE VARIATIONSRECHNUNG.}
\noindent dingungen (5) und (6) auch umgekehrt folgt, dass $\delta\Omega\BolzaEqualAboveGreater0$ für alle zulässigen
Funktionen $\mu$, sodass also aus der ersten Variation keine weiteren Folgerungen
zu ziehen sind. Dagegen hat \textsc{Gauss} \textls[100]{für den speziellen Fall} $U=0$, der
für seine weiteren Schlüsse besonders wichtig ist, mit Hilfe eines in einem
vorangehenden Artikel (Art. 28) von ihm bewiesenen Satzes der Potentialtheorie
gezeigt, dass hier bei der das Minimum liefernden Verteilung \textls[100]{überhaupt
keine unbelegten Stücke der Fläche} vorkommen können, sodass also in
diesem Fall die Gleichung $V=A$ in der ganzen Fläche ohne Ausnahme gilt.
\par\noindent\hspace*{\parindent}5. Auf die sinnreichen Schlüsse, durch die \textsc{Gauss} von diesen Resultaten
aus zum Beweis des im Eingang genannten Existenztheorems der Potentialtheorie gelangt, wollen wir hier nicht eingehen, da sie nicht der Variationsrechnung, sondern ausschliesslich der Potentialtheorie angehören. Dagegen
wollen wir noch auf eine Seite des \textsc{Gauss}schen Variationsproblems hinweisen,
die für die Variationsrechnung von besonderer Bedeutung ist: Die sonst in der
Variationsrechnung behandelten Probleme pflegen auf Differentialgleichungen
als erste notwendige Bedingung des Extremums zu führen. Hier haben wir
es dagegen mit einem Variationsproblem zu tun, das auf eine \textls[100]{Integralgleichung} führt, nämlich auf die Gleichung (5), die ausgeschrieben lautet
\[
\int\frac{m\,ds}{r}=U+A.
\]
\par\noindent\hspace*{\parindent}In der Tat lässt sich das \textsc{Gauss}sche Variationsproblem keinem der üblichen Typen von Variationsproblemen unterordnen, sondern ist ein Beispiel
eines eigenartigen neuen Typus\footnote{Ein ähnliches Problem behandelt D. \textsc{Hilbert} in seinen \textit{Grundzügen einer allgemeinen Theorie
der linearen Differentialgleichungen} (Göttinger Nachrichten 1904, S. 78, 232, 257; 1906, S. 460; 1910, S. 9,
22), und zwar unter dem Namen \textsc{Gauss}'sches Variationsproblem, nämlich die Aufgabe, das Integral
\[
\int_a^b\int_a^b K(s,t)u(s)u(t)\,ds\,dt
\]
zu einem Maximum zu machen mit der Nebenbedingung
\[
\int_a^b(u(s))^2\,ds=1.
\]
}. Denkt man sich nämlich die Koordinaten
eines Punktes der Fläche durch zwei unabhängige Parameter $u$, $v$ ausgedrückt,
die einen bestimmten Bereich $\mathfrak A$ überstreichen, so sieht man, dass das \textsc{Gauss}sche Problem dem folgenden Typus angehört: Unter allen stetigen, nicht
\clearpage\SourceLeaf{540}
\ArithmeticHead{83}{R}{EIN EXISTENZTHEOREM DER POTENTIALTHEORIE.}
\noindent negativen Funktionen $m(u,v)$, die einer isoperimetrischen Bedingung von der
Form
\[
\iint G(u,v)m(u,v)\,du\,dv=M
\]
genügen, diejenige zu bestimmen, die ein vierfaches Integral von der Form
\[
\Omega=\iiiint F(u,v;u',v')m(u,v)m(u',v')\,du\,dv\,du'\,dv'
\]
zu einem Minimum macht, wobei der Integrationsbereich für $u$, $v$ sowohl wie
für $u'$, $v'$ derselbe Bereich $\mathfrak A$ ist. Auch für dieses allgemeinere Problem lässt
sich das \textsc{Gauss}sche Verfahren durchführen. Für die erste Variation erhält man
\[
\begin{aligned}
\delta\Omega={}&\iiiint F(u,v;u',v')\bigl[m(u',v')\mu(u,v)\\
&\qquad+m(u,v)\mu(u',v')\bigr]\,du\,dv\,du'\,dv'.
\end{aligned}
\]
Durch Vertauschung der Bezeichnungen der beiden Reihen von Integrationsvariabeln und gleichzeitige Vertauschung der Integrationsordnungen im zweiten
Summanden rechts, — was der Anwendung des Satzes von Art. 19 beim
\textsc{Gauss}schen Problem entspricht —, lässt sich dieser Ausdruck auf die Form
bringen
\[
\delta\Omega=\iiiint K(u,v;u',v')m(u',v')\mu(u,v)\,du\,dv\,du'\,dv',
\]
wenn man zur Abkürzung
\[
F(u,v;u',v')+F(u',v';u,v)=K(u,v;u',v')
\]
setzt.
\par\noindent\hspace*{\parindent}Eine der \textsc{Gauss}schen Schlussweise ganz analoge Betrachtung führt dann
auf folgende \textls[100]{Integralgleichung erster Art} als erste notwendige Bedingung
des Minimums, wobei $\lambda$ die isoperimetrische Konstante bedeutet,
\[
\iint K(u,v;u',v')m(u',v')\,du'\,dv'=\lambda G(u,v),
\]
die in denjenigen Teilen der Fläche erfüllt sein muss, wo $m>0$, während
das Gleichheitszeichen durch $\BolzaEqualAboveGreater$ zu ersetzen ist in denjenigen Teilen, wo
$m=0$.
\par\noindent\hspace*{\parindent}6. Zusammenfassend können wir also sagen, dass in der \textsc{Gauss}schen
Arbeit die folgenden Punkte für die Variationsrechnung von besonderer
Wichtigkeit sind:
\par\noindent\hspace*{\parindent}1) dass hier bereits die später mit dem Namen \textsc{Dirichlet}sches Prinzip
bezeichnete Schlussweise angewandt wird;
\par\noindent\hfill 11$^{\ast}$\hspace*{5em}\par
\clearpage\SourceLeaf{541}
\ArithmeticHead{84}{L}{BOLZA, GAUSS UND DIE VARIATIONSRECHNUNG.}
\noindent\hspace*{\parindent}2) dass hier ein eigenartiger Typus von Variationsproblemen behandelt
wird, der nicht auf eine Differentialgleichung, sondern auf eine Integralgleichung erster Art als erste notwendige Bedingung eines Extremums führt;
\par\noindent\hspace*{\parindent}3) dass hier wohl zum ersten Mal ein Variationsproblem mit einer Ungleichung als Nebenbedingung vorkommt;
\par\noindent\hspace*{\parindent}4) der Beweis des sogenannten \textsc{du Bois-Reymond}schen Lemmas.
\par\noindent\hspace*{\parindent}Es scheint jedoch nicht, dass die \textsc{Gauss}sche Arbeit, so nachhaltig auch
ihr Einfluss auf die Entwicklung der Potentialtheorie gewesen ist, eine nennenswerte Wirkung auf die Variationsrechnung ausgeübt hat.
\par\begin{center}\rule{.18\linewidth}{.3pt}\end{center}
\clearpage\endgroup
